The output file and directory structure of Corvus is as follows:
Calculations will be run in sub-directories of your current directory named Corvus\_description
where "description" describes what is in the directory. 

Simple calculation that ask directly for a spectrum calculated by FEFF are calculated in the directory
Corvus1\_FEFF, while asking for an average over configurations (target\_list\{ cfavg \}) produces a top-level
directory named,
Corvus3\_cfavg\_spectrum, where spectrum is the named spectrum that was requested by the user, e.g., 
Corvus3\_cfavg\_xanes. Then each of the subdirectories of Corvus3\_cfavg\_xanes, will have directories named by the
index of the absorber as it was run, a label connected to which atom is the absorber, and the code (FEFF in this case). 

For the special target loop, the top level folders start with CorvusN\_loop\_loopparameter\_loopvalue. For example, if the user requests a loop over the parameter feff.corehole, with 
values none, rpa, fsr, then Corvus produces the directories Corvus1\_loop\_feff.corehole\_none, Corvus1\_loop\_feff.corehole\_rpa, Corvus1\_loop\_feff.corehole\_fsr

Finally, the main output files of Corvus are named Corvus.target\_subtarget.out. Thus for calculations that request xanes and provide the structure directly in the input file, the main output
is in Corvus.xanes.out. For configurational averages (cfavg), the main output is in
Corvus.cfavg\_xanes.out (for xanes calculations). Finally for loop calculations, the main output is inside the top-level loop directories, and named the same as it would be for a single calculation, e.g. Corvus.xanes.out or Corvus.cfavg\_xes.out etc. 
